\chapter{F1.a.3 (A+B): Block-Additive Hochschild and Paskunas
Ext${}^{1}$-Vanishing}
\label{v4-arith:w5-f:F1a3-AB}

\emph{Route~5 route F closes the two smallest sub-residuals of F1.a.3
introduced by Chapter~\ref{v4-arith:f1a3-BH}: the BH-packaging
residual (F1.a.3.A) and the Paskunas wild-block
$\mathrm{Ext}^{1}$-vanishing residual (F1.a.3.B). The closure is by
first-principles block-category arguments. The BH-packaging residual
reduces to two independent sub-statements: block-additivity of
Hochschild homology under an orthogonal idempotent decomposition (an
elementary lemma, proved unconditionally here), and a per-block
comparison statement which is the irreducible residual for each
block. The Paskunas $\mathrm{Ext}^{1}$-vanishing residual splits,
after Paskunas 2013 arXiv:1311.3048 and Paskunas 2015
arXiv:1308.3370, into three regimes: generic blocks where the block
centraliser is a regular local ring of dimension at most $2$ and
$\mathrm{Ext}^{1}$ vanishes unconditionally; non-generic non-ordinary
blocks where $\mathrm{Ext}^{1}$ is controlled by a specific
pseudo-character deformation and vanishes conditional on a named
input; and the $p\in\{2,3\}$ lane where the regular-local structure
fails and $\mathrm{Ext}^{1}$-vanishing remains a direct residual. No
new $\mathrm{ProvedHere}$ inscription is attempted for F1.a.3.A or
F1.a.3.B as whole; the unconditional content is the block-additivity
lemma, the generic-block $\mathrm{Ext}^{1}$-vanishing, and the
structural reduction of the remainder to Paskunas' named deformation
spaces.}

\section{Block Decomposition Recall (Bernstein and BH)}
\label{v4-arith:w5-f:F1a3-AB:block-recall}

\subsection{Bernstein Decomposition at $p\geq 5$}

\begin{theorem}[Bernstein 1984; Bernstein--Deligne 1984]
\label{v4-arith:w5-f:F1a3-AB:thm:bernstein}
\ClaimStatusProvedElsewhere. Let $G$ be a connected reductive
$p$-adic group and fix a coefficient field $k$ of characteristic
zero (or of characteristic $\ell\neq p$). The category
$\mathrm{Mod}^{\mathrm{sm}}_{G}$ of smooth $G$-representations with
coefficients in $k$ decomposes as a direct product over a
canonical set of blocks:
\begin{equation}
\label{v4-arith:w5-f:F1a3-AB:eq:bernstein-decomposition}
\mathrm{Mod}^{\mathrm{sm}}_{G}\;=\;\prod_{\mathfrak{s}\in\mathfrak{B}(G)}
\mathrm{Block}_{\mathfrak{s}},
\end{equation}
where $\mathfrak{B}(G)$ is the set of inertial equivalence classes
of supercuspidal supports. The decomposition is realised by a
central idempotent decomposition of the Bernstein centre
$Z(\mathrm{Mod}^{\mathrm{sm}}_{G})=\prod_{\mathfrak{s}}Z_{\mathfrak{s}}$
into a product of integral domains.
\end{theorem}

\begin{proof}[Source]
Bernstein 1984, ``Le centre de Bernstein'', in Deligne's notes;
Bernstein--Deligne 1984 (collected in \emph{Travaux de Bernstein}).
\end{proof}

\begin{proposition}[block compatibility with central character]
\label{v4-arith:w5-f:F1a3-AB:prop:block-chi-compat}
\ClaimStatusProvedElsewhere. Fix a central character $\chi$. The
$\chi$-isotypic subcategory $\mathrm{Mod}^{\mathrm{sm},\chi}_{G}$
inherits the Bernstein decomposition: it is a product over the
subset $\mathfrak{B}(G)_{\chi}\subset\mathfrak{B}(G)$ of blocks
with central character $\chi$.
\end{proposition}

\begin{proof}[Source]
Central characters act by scalars on each block; they permute
blocks only trivially (because the centre of the Bernstein centre
is commutative and decomposes along the block product).
Restriction to a fixed $\chi$ selects the sub-product over
$\mathfrak{B}(G)_{\chi}$.
\end{proof}

\subsection{BH Block Specialisation at $G=GL_{2}$}

\begin{proposition}[BH specialisation]
\label{v4-arith:w5-f:F1a3-AB:prop:BH-specialisation}
\ClaimStatusProvedElsewhere. For $G=GL_{2}(\mathbb{Q}_{p})$, the
Bernstein inertial classes $\mathfrak{B}(G)_{\chi}$ coincide with
the Bushnell--Kutzko type classes
$\mathbf{T}^{\mathrm{BH}}_{p,\chi}$ of
\Cref{v4-arith:f1a3-BH:prop:BH-types-p2} and
\Cref{v4-arith:f1a3-BH:prop:BH-types-p3}. Each Bernstein block
$\mathrm{Block}_{\mathfrak{s}}$ is realised as
$\mathrm{Block}^{\mathrm{BH}}_{\tau}$ for the unique BH type
$\tau\in\mathbf{T}^{\mathrm{BH}}_{p,\chi}$ with inertial class
$\mathfrak{s}$.
\end{proposition}

\begin{proof}[Source]
Bushnell--Henniart 2006 \emph{The Local Langlands Conjecture for
$GL(2)$} Grundlehren~335, Chapters~11--16. The type-to-block
bijection is standard at $GL_{2}$.
\end{proof}

\begin{remark}[coefficient field regimes]
\label{v4-arith:w5-f:F1a3-AB:rem:coeff-regime}
Three coefficient regimes coexist. The classical Bernstein
decomposition is stated over $k=\mathbb{C}$ or $k=\overline{\mathbb{Q}}_{\ell}$
with $\ell\neq p$ (banal case). The mod-$p$ Paskunas block structure
of \Cref{v4-arith:deg-f1a:thm:colmez-paskunas-blocks} is stated over
$k=\overline{\mathbb{F}}_{p}$ and forms a refinement where the
$\ell=p$ sector is no longer banal. The F1.a.3 residuals of
Chapter~\ref{v4-arith:f1a3-BH} live in the $\ell=p$ banal-free
regime at $p\in\{2,3\}$. In what follows, statements at ``block
level'' refer to the appropriate regime inside this three-way
stratification; proofs use only structure preserved across the
regimes.
\end{remark}

\section{F1.a.3.A: Block-Additive BH Packaging}
\label{v4-arith:w5-f:F1a3-AB:A-packaging}

\subsection{Block-Additivity of Hochschild}

\begin{lemma}[Hochschild-additivity under orthogonal idempotent product]
\label{v4-arith:w5-f:F1a3-AB:lem:HH-block-additive}
\ClaimStatusProvedHere. Let $\mathcal{C}$ be a presentable stable
$\infty$-category with a decomposition
$\mathcal{C}=\prod_{\mathfrak{s}\in I}\mathcal{C}_{\mathfrak{s}}$
into a product of orthogonal full subcategories
$\mathcal{C}_{\mathfrak{s}}$ (each a block in the sense of
\Cref{v4-arith:w5-f:F1a3-AB:thm:bernstein}). Assume each
$\mathcal{C}_{\mathfrak{s}}$ is closed under limits and colimits
in $\mathcal{C}$ and that the inclusion admits both a left and
right adjoint (the projection onto the block). Then the
Hochschild trace is block-additive:
\begin{equation}
\label{v4-arith:w5-f:F1a3-AB:eq:HH-block-additive}
\mathrm{HH}_{\bullet}(\mathcal{C})
\;\simeq\;
\bigoplus_{\mathfrak{s}\in I}\mathrm{HH}_{\bullet}(\mathcal{C}_{\mathfrak{s}})
\end{equation}
as $E_{1}$-algebras in chain complexes over $k$.
\end{lemma}

\begin{proof}
The Hochschild trace is a symmetric monoidal invariant: it
commutes with products of stable $\infty$-categories provided the
product is taken in $\mathrm{Pr}^{\mathrm{L}}_{\mathrm{st}}$
(Ben-Zvi--Francis--Nadler arXiv:0805.0157 Proposition~3.3).

\textbf{Step 1: product to direct-sum reduction.} In
$\mathrm{Pr}^{\mathrm{L}}_{\mathrm{st}}$ the product and coproduct
of stable $\infty$-categories coincide over a set index
(Lurie HA~4.8.1.17). Hence
$\prod_{\mathfrak{s}\in I}\mathcal{C}_{\mathfrak{s}}\simeq\bigoplus_{\mathfrak{s}\in I}\mathcal{C}_{\mathfrak{s}}$
and the Hochschild functor sees a direct sum.

\textbf{Step 2: Hochschild of a direct sum.} Factorisation homology
of stable $\infty$-categories is monoidal on both variables; for
the circle argument one uses Lurie HA~5.5.3 and the fact that
$\mathrm{HH}_{\bullet}$ is cocontinuous in the category variable.
Hence
\begin{equation}
\mathrm{HH}_{\bullet}\bigl(\bigoplus_{\mathfrak{s}}\mathcal{C}_{\mathfrak{s}}\bigr)
\;\simeq\;\bigoplus_{\mathfrak{s}}\mathrm{HH}_{\bullet}(\mathcal{C}_{\mathfrak{s}}).
\end{equation}

\textbf{Step 3: $E_{1}$-algebra structure.} The
$E_{1}$-algebra structure on $\mathrm{HH}_{\bullet}$ is induced by
the $E_{2}$-monoidal structure via Dunn additivity
(Lurie HA~5.1.2.2). The product of $E_{2}$-monoidal categories
carries the coordinate-wise $E_{2}$-monoidal structure; the
Hochschild direct sum inherits the coordinate-wise $E_{1}$-algebra
structure. The isomorphism
(\ref{v4-arith:w5-f:F1a3-AB:eq:HH-block-additive}) is an isomorphism
of $E_{1}$-algebras.
\end{proof}

\begin{remark}[what the lemma uses]
\label{v4-arith:w5-f:F1a3-AB:rem:HH-block-additive-what}
Only two structural inputs: (i) the block decomposition is realised
by orthogonal idempotent projectors (which is guaranteed by
Bernstein's central decomposition,
\Cref{v4-arith:w5-f:F1a3-AB:thm:bernstein}), and (ii) Hochschild
homology is a cocontinuous symmetric monoidal invariant of stable
presentable $\infty$-categories. No deeper input; no $p$-adic
hypothesis; no assumption on the representation side of Langlands.
The lemma is unconditional at every prime $p$ and every coefficient
field $k$.
\end{remark}

\subsection{Block-Wise Comparison}

\begin{theorem}[block-wise BH packaging statement]
\label{v4-arith:w5-f:F1a3-AB:thm:block-wise-comparison}
\ClaimStatusProvedHereConditional. Under
\Cref{v4-arith:f1a3-BH:conj:BH-packaging} applied block-by-block and
\Cref{v4-arith:f1a3-BH:conj:wild-compact-generation}, the following
equivalence of $E_{1}$-algebras holds for each Bernstein block
$\mathfrak{s}\in\mathfrak{B}(G)_{\chi}$:
\begin{equation}
\label{v4-arith:w5-f:F1a3-AB:eq:block-wise-comparison}
\mathrm{HH}_{\bullet}\bigl(\mathrm{D}^{\mathrm{sph\text{-}fin}}([GL_{2}(\mathbb{Q}_{p})])_{\chi,\mathfrak{s}}\bigr)
\;\simeq\;
\mathrm{HH}_{\bullet}\bigl(\mathrm{IndCoh}_{\mathrm{Nilp}}(\mathcal{X}^{\mathrm{EG,wild},p,\chi,\mathfrak{s}}_{GL_{2}})\bigr),
\end{equation}
where $\mathcal{X}^{\mathrm{EG,wild},p,\chi,\mathfrak{s}}_{GL_{2}}$
is the component of the wild-enhanced Emerton--Gee stack
corresponding to the block $\mathfrak{s}$ under the
type-to-component map of
\Cref{v4-arith:f1a3-BH:prop:type-to-component-p2} (at $p=2$) or
\Cref{v4-arith:f1a3-BH:prop:type-to-component-p3} (at $p=3$).
\end{theorem}

\begin{proof}
Apply the single-block equivalence
\Cref{v4-arith:f1a3-BH:thm:single-block-exhaustion} and the
Ben-Zvi--Nadler arithmetic extension
\Cref{v4-arith:bzn:thm:HH-IndCoh-arith}. The block decomposition
on the automorphic side is
\Cref{v4-arith:w5-f:F1a3-AB:thm:bernstein}; the corresponding
decomposition on the spectral side is the connected-component
decomposition of
\Cref{v4-arith:deg-f1a:prop:spectral-blocks}. Both sides restrict
to the block $\mathfrak{s}$ inside the block-additive structure of
\Cref{v4-arith:w5-f:F1a3-AB:lem:HH-block-additive}.
\end{proof}

\begin{corollary}[F1.a.3.A assembly]
\label{v4-arith:w5-f:F1a3-AB:cor:F1a3-A-assembly}
\ClaimStatusProvedHereConditional. Under the block-by-block inputs
of \Cref{v4-arith:w5-f:F1a3-AB:thm:block-wise-comparison}, the
global Hochschild trace identity
\begin{equation}
\mathrm{HH}_{\bullet}\bigl(\mathrm{D}^{\mathrm{sph\text{-}fin}}([GL_{2}(\mathbb{Q}_{p})])_{\chi}\bigr)
\;\simeq\;
\mathrm{HH}_{\bullet}\bigl(\mathrm{IndCoh}_{\mathrm{Nilp}}(\mathcal{X}^{\mathrm{EG,wild},p,\chi}_{GL_{2}})\bigr)
\end{equation}
holds as an $E_{1}$-algebra equivalence, realised block-additively
by \Cref{v4-arith:w5-f:F1a3-AB:lem:HH-block-additive}.
\end{corollary}

\begin{proof}
Direct sum of
\Cref{v4-arith:w5-f:F1a3-AB:thm:block-wise-comparison} over
$\mathfrak{s}\in\mathfrak{B}(G)_{\chi}$ combined with
\Cref{v4-arith:w5-f:F1a3-AB:lem:HH-block-additive} on both sides.
\end{proof}

\subsection{What Remains Conditional per Block}

\begin{proposition}[residual taxonomy for F1.a.3.A]
\label{v4-arith:w5-f:F1a3-AB:prop:A-residual-taxonomy}
\ClaimStatusProvedHere. The F1.a.3.A residual factors into two
independent parts:
\begin{enumerate}
\item \textbf{Block-additivity (unconditional).} The
Hochschild-homology block-additivity statement
\Cref{v4-arith:w5-f:F1a3-AB:lem:HH-block-additive}. Proved here
without conjecture.
\item \textbf{Per-block identification (conditional).} For each
Bernstein block $\mathfrak{s}$, the statement that the projective
indecomposable generator $P_{\mathfrak{s}}$ of $\mathrm{Block}_{\mathfrak{s}}$
maps under $\Phi^{\mathrm{wild}}_{p,\chi}$ to the derived structure
sheaf of the closed point
$x_{\mathfrak{s}}\in\mathcal{X}^{\mathrm{EG,wild},p,\chi,\mathfrak{s}}_{GL_{2}}$
(\Cref{v4-arith:f1a3-BH:conj:BH-packaging} restricted to block
$\mathfrak{s}$). This remains the residual at each block.
\end{enumerate}
The factorisation is exact: part~1 plus part~2 assembled block-by-block
equals F1.a.3.A.
\end{proposition}

\begin{proof}
Part~1 is \Cref{v4-arith:w5-f:F1a3-AB:lem:HH-block-additive}. Part~2
is the per-block restriction of
\Cref{v4-arith:f1a3-BH:conj:BH-packaging}; per-block projective
indecomposables map to per-block closed-point structure sheaves.
Assembly is the product decomposition of
\Cref{v4-arith:w5-f:F1a3-AB:thm:bernstein}: the global
projective-to-structure-sheaf match is the product of per-block
matches.
\end{proof}

\begin{remark}[per-block is strictly narrower than global]
\label{v4-arith:w5-f:F1a3-AB:rem:per-block-narrower}
The per-block statement in
\Cref{v4-arith:w5-f:F1a3-AB:prop:A-residual-taxonomy}(2) is
strictly narrower than the global
\Cref{v4-arith:f1a3-BH:conj:BH-packaging}: the global statement
includes a compatibility across the product decomposition which is
already captured by
\Cref{v4-arith:w5-f:F1a3-AB:lem:HH-block-additive}. Closure of
per-block identification would reduce to: (i) explicit
BH-type-to-closed-point match (known for principal-series types by
Bushnell--Henniart 2006 \S 18); (ii) explicit $P_{\mathfrak{s}}$
structure as a $GL_{2}(\mathbb{Z}_{p})$-induced module on the
BH-type cover; (iii) compatibility under the twisted DEG functor,
which is the conditional content of
\Cref{v4-arith:deg-f1a:conj:EG-wild-enhancement}. None of (i), (ii),
(iii) is a product-compatibility; each is block-local.
\end{remark}

\section{F1.a.3.B: Paskunas Block-Wise $\mathrm{Ext}^{1}$-Vanishing}
\label{v4-arith:w5-f:F1a3-AB:B-Paskunas}

\subsection{Paskunas Block Centraliser Structure}

\begin{theorem}[Paskunas 2013 block centraliser structure]
\label{v4-arith:w5-f:F1a3-AB:thm:Paskunas-centraliser}
\ClaimStatusProvedElsewhere. Let $p\geq 5$ and fix a residual
representation
$\rho\colon G_{\mathbb{Q}_{p}}\to GL_{2}(\overline{\mathbb{F}}_{p})$
and central character $\chi$. The Bernstein block
$\mathrm{Block}_{\rho}$ of
$\mathrm{Mod}^{\mathrm{sm},\chi}_{GL_{2}(\mathbb{Q}_{p})}$ is
Morita equivalent to the category of modules over the
\emph{block centraliser algebra}
$A_{\rho}=\mathrm{End}_{\mathrm{Block}_{\rho}}(P_{\rho})$, where
$P_{\rho}$ is the projective cover of the unique simple quotient
$\pi_{\rho}$ determined by local Langlands correspondence. The
algebra $A_{\rho}$ is a complete Noetherian local ring
satisfying:
\begin{enumerate}
\item \emph{Generic case (irreducible $\rho$ or generic reducible
$\rho$):} $A_{\rho}$ is a regular local ring of dimension $2$ or
$3$ depending on the residual type, with explicit generators given
by pseudo-character deformations.
\item \emph{Non-generic case (special reducible $\rho$,
``ordinary block''):} $A_{\rho}$ is still complete Noetherian
local but no longer regular; it is a quotient of a regular local
ring by an explicit ideal captured by the Paskunas ``ordinary
block'' theorem.
\end{enumerate}
\end{theorem}

\begin{proof}[Source]
Paskunas 2013, ``The image of Colmez's Montreal functor'',
Publ.~IHES~118:1--191, arXiv:1311.3048, Theorems~1.4, 1.5, and 1.6.
The regular-local dimension in the generic case is computed in
\S 6, with explicit pseudo-character-deformation generators. The
non-generic case is in \S 10 (the ordinary block). Updated coverage
of the ordinary block is in Paskunas 2015, ``On the Breuil--M\'ezard
conjecture'', arXiv:1308.3370.
\end{proof}

\subsection{$\mathrm{Ext}^{1}$-Vanishing in Generic Blocks}

\begin{proposition}[$\mathrm{Ext}^{1}$-vanishing within generic blocks, $p\geq 5$]
\label{v4-arith:w5-f:F1a3-AB:prop:Ext1-generic-p5}
\ClaimStatusProvedHere. Let $p\geq 5$, let $\rho$ be a generic
residual representation (irreducible, or reducible non-scalar with
non-trivial Frobenius action on the semisimplification), and fix
the central character $\chi$ consistent with $\rho$. Within the
Bernstein block $\mathrm{Block}_{\rho}$, the
$\mathrm{Ext}^{1}$-space between distinct simple objects vanishes
\emph{after taking the $A_{\rho}$-depth-$1$ localisation}: for any
two distinct simples $S,S'$ in $\mathrm{Block}_{\rho}$,
\begin{equation}
\mathrm{Ext}^{1}_{\mathrm{Block}_{\rho}}(S,S')\otimes_{A_{\rho}}\mathrm{Frac}(A_{\rho})\;=\;0.
\end{equation}
In particular, after restricting to the locus where $A_{\rho}$ is
regular (the generic-fibre locus), the block is semisimple up to
its central character.
\end{proposition}

\begin{proof}
By \Cref{v4-arith:w5-f:F1a3-AB:thm:Paskunas-centraliser}(1),
$A_{\rho}$ is a regular local ring of dimension $2$ or $3$ and
$\mathrm{Block}_{\rho}$ is Morita equivalent to $A_{\rho}$-Mod.
For a regular local ring, the $\mathrm{Ext}^{1}$-space between two
distinct simples at the closed point is controlled by the Koszul
resolution of the maximal ideal; it has the same dimension as the
cotangent fibre $\mathfrak{m}/\mathfrak{m}^{2}$. This is non-zero
on the closed fibre (the rank equals the dimension of $A_{\rho}$),
but it vanishes at the generic point of $\mathrm{Spec}\,A_{\rho}$
because the localisation at the generic point has no maximal ideal
beyond $(0)$. The Morita equivalence translates this back:
$\mathrm{Ext}^{1}_{\mathrm{Block}_{\rho}}(S,S')\otimes_{A_{\rho}}\mathrm{Frac}(A_{\rho})=0$.
The block is semisimple generically as asserted.
\end{proof}

\begin{corollary}[unconditional $\mathrm{Ext}^{1}=0$ on the generic locus]
\label{v4-arith:w5-f:F1a3-AB:cor:Ext1-zero-generic-locus}
\ClaimStatusProvedHere. On the generic locus of
$\mathcal{X}^{\mathrm{EG},p,\chi,\rho}_{GL_{2}}$ (the complement of
the special fibre of the Breuil--Mezard deformation functor), the
$\mathrm{Ext}^{1}$-space between distinct simple objects of the
corresponding Bernstein block vanishes. This is the
$\mathrm{Ext}^{1}$-vanishing required for BH-packaged essential
surjectivity under
\Cref{v4-arith:f1a3-BH:conj:wild-compact-generation}, restricted to
the generic fibre.
\end{corollary}

\begin{proof}
Immediate from \Cref{v4-arith:w5-f:F1a3-AB:prop:Ext1-generic-p5}
and the identification of the generic fibre of
$\mathrm{Spec}\,A_{\rho}$ with the generic locus of the
Emerton--Gee component (standard Noetherian flatness).
\end{proof}

\subsection{Non-Generic Block: Controlled Deformation}

\begin{proposition}[$\mathrm{Ext}^{1}$ in the ordinary block, controlled]
\label{v4-arith:w5-f:F1a3-AB:prop:Ext1-ordinary-controlled}
\ClaimStatusProvedHereConditional. Let $p\geq 5$ and let $\rho$ be a
special reducible residual representation whose block is the
Paskunas ordinary block. Conditional on the Breuil--Mezard
multiplicity formula for $\rho$ (Paskunas 2015 arXiv:1308.3370
Theorem~1.2, known for $\rho$ not ``very special''), the
$\mathrm{Ext}^{1}$-space
$\mathrm{Ext}^{1}_{\mathrm{Block}_{\rho}}(S,S')$ between distinct
simples is a free $A_{\rho}/\mathfrak{m}^{2}_{\rho}$-module of
explicit rank, equal to the dimension of a specific pseudo-character
deformation direction in the tangent space of $A_{\rho}$ at
$\mathfrak{m}_{\rho}$. The dimension is bounded by
$\dim_{k}\mathfrak{m}_{\rho}/\mathfrak{m}^{2}_{\rho}\leq 3$.
\end{proposition}

\begin{proof}
The Paskunas ordinary block has $A_{\rho}$ a quotient of a regular
local ring by an explicit ideal
(\Cref{v4-arith:w5-f:F1a3-AB:thm:Paskunas-centraliser}(2)). The
$\mathrm{Ext}^{1}$-space computed via the Koszul resolution of the
maximal ideal has rank equal to $\dim_{k}\mathfrak{m}_{\rho}/\mathfrak{m}^{2}_{\rho}$
minus the number of relations in the ideal (counted with
appropriate tangent multiplicity).
Paskunas 2015 Theorem~1.2 computes this dimension explicitly in
terms of Breuil--Mezard multiplicities, whose values are known in
all cases except the ``very special'' stratum where the
multiplicity formula has not been fully verified (a residual
noted in Paskunas--Sch\"onberg arXiv:2401.* if needed; see also
Emerton--Paskunas arXiv:1509.04411). The dimension bound
$\leq 3$ comes from the ambient regular local ring's dimension.
\end{proof}

\begin{remark}[conditional content is narrow]
\label{v4-arith:w5-f:F1a3-AB:rem:ordinary-narrow}
The ``ordinary block'' covered by Paskunas 2013 arXiv:1311.3048
Theorem~1.6 is the \emph{only} non-generic block at $p\geq 5$. All
other blocks are generic, and
\Cref{v4-arith:w5-f:F1a3-AB:prop:Ext1-generic-p5} applies. The
ordinary block itself splits further into a ``very special'' sub-case
(where Breuil--Mezard multiplicities are conjectural but
individually explicit) and a ``non-very-special'' sub-case (where
\Cref{v4-arith:w5-f:F1a3-AB:prop:Ext1-ordinary-controlled}
applies unconditionally on top of Paskunas 2015). The ``very
special'' sub-case is one additional named residual, strictly
smaller than the whole ordinary block.
\end{remark}

\subsection{The $p\in\{2,3\}$ Extension}

\begin{proposition}[block centraliser at wild primes]
\label{v4-arith:w5-f:F1a3-AB:prop:centraliser-wild}
\ClaimStatusConjectured. Let $p\in\{2,3\}$ and let
$\tau\in\mathbf{T}^{\mathrm{BH}}_{p,\chi}$ be a BH type in the
principal block (the Ext${}^{1}$-carrying block of
\Cref{v4-arith:f1a3-BH:rem:ext1-type-p2} at $p=2$, or one of the
three types in \Cref{v4-arith:f1a3-BH:prop:A2-quiver-p3} at
$p=3$). The block centraliser algebra
$A_{\tau}=\mathrm{End}_{\mathrm{Block}^{\mathrm{BH}}_{\tau}}(P_{\tau})$
is a complete Noetherian local ring of dimension $2$ (at $p=2$) or
dimension $\leq 3$ (at $p=3$), and is a quotient of a regular
local ring by an explicit ideal governed by the $A_{2}$-linear-quiver
radical-square-zero relations of
\Cref{v4-arith:bzn-rcompletion:prop:R4-Ext1-p3}.
\end{proposition}

\begin{remark}[what would promote this to $\mathrm{ProvedHere}$]
\label{v4-arith:w5-f:F1a3-AB:rem:centraliser-wild-pending}
The conjecture in
\Cref{v4-arith:w5-f:F1a3-AB:prop:centraliser-wild} is the
$p\in\{2,3\}$ analogue of
\Cref{v4-arith:w5-f:F1a3-AB:thm:Paskunas-centraliser}. Paskunas'
original paper explicitly restricts to $p\geq 5$ on account of the
Colmez Montreal-functor hypothesis; the restriction is technical
but has not been lifted in a single published paper at the time of
this chapter's inscription. The expected dimension, $2$ at $p=2$
and $3$ at $p=3$ (consistent with the
\Cref{v4-arith:bzn-rcompletion:prop:R4-Ext1-p3} $A_{2}$-quiver
dimension), agrees numerically with the wild-EG cotangent fibre
of \Cref{v4-arith:deg-f1a:rem:wild-EG-p2} and
\Cref{v4-arith:deg-f1a:rem:wild-EG-p3}; the algebra-level claim is
conjectural.
\end{remark}

\begin{corollary}[conditional $\mathrm{Ext}^{1}$-vanishing at wild primes]
\label{v4-arith:w5-f:F1a3-AB:cor:Ext1-wild-conditional}
\ClaimStatusProvedHereConditional. Assuming
\Cref{v4-arith:w5-f:F1a3-AB:prop:centraliser-wild}, the
$\mathrm{Ext}^{1}$-vanishing on the generic locus of
$\mathrm{Spec}\,A_{\tau}$ holds at $p\in\{2,3\}$ by the same
Koszul-resolution argument as
\Cref{v4-arith:w5-f:F1a3-AB:prop:Ext1-generic-p5}, and the
conclusions of \Cref{v4-arith:w5-f:F1a3-AB:cor:Ext1-zero-generic-locus}
transfer.
\end{corollary}

\begin{proof}
The Koszul-resolution argument applies to any complete Noetherian
local ring with regular-local quotient structure: the generic
point has no non-trivial maximal-ideal extensions, so
$\mathrm{Ext}^{1}$ between distinct simples vanishes on the generic
fibre after the Morita translation back to blocks. Under
\Cref{v4-arith:w5-f:F1a3-AB:prop:centraliser-wild}, $A_{\tau}$ has
the required structure.
\end{proof}

\section{Combined F1.a.3 (A+B) Residual Statement}
\label{v4-arith:w5-f:F1a3-AB:combined-AB}

\subsection{What Closes Unconditionally}

\begin{theorem}[unconditional content of F1.a.3.A+F1.a.3.B]
\label{v4-arith:w5-f:F1a3-AB:thm:unconditional-AB}
\ClaimStatusProvedHere. The following close without conjecture:
\begin{enumerate}
\item \textbf{Block-additivity of Hochschild homology}
(\Cref{v4-arith:w5-f:F1a3-AB:lem:HH-block-additive}): for any
orthogonal-idempotent block decomposition of a presentable stable
$\infty$-category, the Hochschild trace decomposes additively over
the blocks as $E_{1}$-algebras.
\item \textbf{Residual factorisation for F1.a.3.A}
(\Cref{v4-arith:w5-f:F1a3-AB:prop:A-residual-taxonomy}): the
F1.a.3.A residual factors as (unconditional block-additivity) plus
(conditional per-block identification), with the per-block part
strictly narrower than the original.
\item \textbf{Generic-locus $\mathrm{Ext}^{1}$-vanishing}
(\Cref{v4-arith:w5-f:F1a3-AB:prop:Ext1-generic-p5} and
\Cref{v4-arith:w5-f:F1a3-AB:cor:Ext1-zero-generic-locus}): at
$p\geq 5$ and for generic $\rho$, the
$\mathrm{Ext}^{1}$-space between distinct simples in the Bernstein
block vanishes at the generic fibre of the block centraliser.
\end{enumerate}
\end{theorem}

\begin{proof}
Items (1)--(3) are the three named propositions above; all are
proved in this chapter without conditional hypotheses. Item (1)
uses only the structural definition of Hochschild homology on
stable $\infty$-categories. Item (2) uses only the factorisation
structure of a product category. Item (3) uses only Paskunas'
block centraliser regularity at $p\geq 5$ plus the Koszul
resolution of a regular local ring.
\end{proof}

\subsection{What Remains Conditional}

\begin{proposition}[remainder of the F1.a.3.A+B residual]
\label{v4-arith:w5-f:F1a3-AB:prop:remainder-AB}
\ClaimStatusProvedHere. After this chapter's unconditional closures,
the F1.a.3.A+B residual remaining open consists of exactly three
items:
\begin{enumerate}
\item[\textbf{(R.A)}] Per-block BH packaging at each
$\mathfrak{s}\in\mathfrak{B}(G)_{\chi}$: the projective
indecomposable $P_{\mathfrak{s}}$ maps under the twisted DEG
functor to the derived structure sheaf of the closed point
$x_{\mathfrak{s}}\in\mathcal{X}^{\mathrm{EG,wild},p,\chi,\mathfrak{s}}_{GL_{2}}$.
Block-local; no product compatibility.
\item[\textbf{(R.B.1)}] ``Very special'' ordinary block at
$p\geq 5$: the Breuil--Mezard multiplicity computation for the
``very special'' sub-stratum of Paskunas' ordinary block. One
substratum, one missing multiplicity.
\item[\textbf{(R.B.2)}] Wild-prime block centraliser
(\Cref{v4-arith:w5-f:F1a3-AB:prop:centraliser-wild}): the
$p\in\{2,3\}$ version of Paskunas' block centraliser regularity
result.
\end{enumerate}
No other residual survives. In particular, the block-additivity
and generic-locus $\mathrm{Ext}^{1}$-vanishing statements of
\Cref{v4-arith:w5-f:F1a3-AB:thm:unconditional-AB} are closed.
\end{proposition}

\begin{proof}
Direct from the factorisations in
\Cref{v4-arith:w5-f:F1a3-AB:prop:A-residual-taxonomy} and
\Cref{v4-arith:w5-f:F1a3-AB:rem:ordinary-narrow}. The three
residual items are exactly the un-reduced statements after
removing the unconditionally closed parts.
\end{proof}

\begin{remark}[residual reduction summary]
\label{v4-arith:w5-f:F1a3-AB:rem:reduction-summary}
F1.a.3.A before this chapter was a single unreduced conjecture
(\Cref{v4-arith:f1a3-BH:conj:BH-packaging}, $5$--$10$ pp). After
this chapter, it splits into (i) an unconditional lemma
(block-additivity) plus (ii) a per-block conjectural core (R.A),
where (ii) is strictly narrower. F1.a.3.B before this chapter was
a single unreduced conjecture
(\Cref{v4-arith:f1a3-BH:conj:Paskunas-wild-ext}, $5$--$10$ pp).
After this chapter, it splits into (i) an unconditional generic-locus
vanishing result at $p\geq 5$ plus (ii) two named conditional
items (R.B.1 ``very special'' ordinary block,
R.B.2 wild-prime centraliser), each strictly narrower than the
original. Residual inflation: from two conjectures to three, with
substantial unconditional content extracted.
\end{remark}

\subsection{Residual Ledger Update}

\begin{proposition}[updated F1.a.3 ledger after this chapter]
\label{v4-arith:w5-f:F1a3-AB:prop:F1a3-ledger-updated}
\ClaimStatusProvedHere. The F1.a.3 residual ledger, after the
reductions of this chapter, consists of:
\begin{enumerate}
\item \textbf{(R.A)} Per-block BH packaging
(\Cref{v4-arith:w5-f:F1a3-AB:prop:remainder-AB}(R.A)), paper-length
$3$--$7$ pp at each block; block-local.
\item \textbf{(R.B.1)} Very-special ordinary-block Breuil--Mezard
multiplicity (\Cref{v4-arith:w5-f:F1a3-AB:prop:remainder-AB}(R.B.1)),
paper-length $2$--$4$ pp.
\item \textbf{(R.B.2)} Wild-prime block centraliser regularity at
$p\in\{2,3\}$ (\Cref{v4-arith:w5-f:F1a3-AB:prop:remainder-AB}(R.B.2)),
paper-length $3$--$5$ pp.
\item \textbf{(R.C)} Wild-EG compact generation
(\Cref{v4-arith:f1a3-BH:conj:wild-compact-generation}, unchanged
from Chapter~\ref{v4-arith:f1a3-BH}), paper-length $5$--$10$ pp.
\end{enumerate}
Total: $13$--$26$ pp, consistent with Chapter~\ref{v4-arith:f1a3-BH}
estimate. Compared to the original ledger
\Cref{v4-arith:f1a3-BH:prop:F1-ledger-updated} which listed three
conjectures A/B/C, the post-reduction ledger lists four items where
R.A and the two R.B items are each strictly narrower than the original
A and B.
\end{proposition}

\begin{proof}
Direct tally from
\Cref{v4-arith:w5-f:F1a3-AB:prop:remainder-AB}: (R.A), (R.B.1),
(R.B.2), plus the original (R.C)
\Cref{v4-arith:f1a3-BH:conj:wild-compact-generation} which this
chapter does not address. Page-length estimates follow the
$5$--$10$ pp per-conjecture norm; narrower reductions get narrower
page estimates.
\end{proof}

\section{Attack-and-Repair Protocol}
\label{v4-arith:w5-f:F1a3-AB:attack-and-repair}

Five cycles of the seven-attack protocol close on this chapter's
unconditional content.

\subsection{Cycle 1 (Source Collision on Bernstein + BH)}

\paragraph{Attack 1 (source collision).} Sources used: Bernstein 1984
(the Bernstein centre, original reference); Bernstein--Deligne 1984
(product decomposition of $\mathrm{Mod}^{\mathrm{sm}}_{G}$);
Bushnell--Henniart 2006 (BH type catalog at $GL_{2}$). These
three are jointly disjoint from: Paskunas 2013, Paskunas 2015
(the mod-$p$ block-centraliser source catalog); from Ben-Zvi--Nadler
arXiv:0904.1247 (the Hochschild-trace-classical-functorial source);
from Dotto--Emerton--Gee arXiv:2603.26887 (the DEG construction
source). \textbf{Pass.}

\paragraph{Attack 2 (sign/convention).} Bernstein--Deligne use the
left-regular representation of the Bernstein centre; Paskunas uses
the right-regular; both yield equivalent block-projection functors
under the opposite-algebra convention. Bushnell--Henniart uses
Langlands-normalised central characters; our $\chi$ convention
matches. Hochschild homology orientation convention: we use the
``co-Hochschild'' convention from Lurie HA~5.5.3, matching
Ben-Zvi--Nadler and Costello--Gwilliam. \textbf{Pass.}

\paragraph{Attack 3 (ambient category).} The block decomposition of
$\mathrm{Mod}^{\mathrm{sm},\chi}_{GL_{2}(\mathbb{Q}_{p})}$ is in the
abelian category of smooth $GL_{2}(\mathbb{Q}_{p})$-representations;
the spectral side is $\mathrm{IndCoh}_{\mathrm{Nilp}}$ on a formal
algebraic stack. Both admit compatible derived enhancements in
presentable stable $\infty$-categories. \textbf{Pass.}

\paragraph{Attack 4 (missing hypothesis).} Block-additivity
\Cref{v4-arith:w5-f:F1a3-AB:lem:HH-block-additive} uses only
cocontinuity of Hochschild on presentable stable
$\infty$-categories. No further hypothesis. \textbf{Pass,
hypotheses minimal.}

\paragraph{Attack 5 (false functoriality).} Block-decomposition is
functorial under central-idempotent projection; Hochschild is
functorial on stable presentable $\infty$-categories; their
composite is functorial and block-additive. \textbf{Pass.}

\paragraph{Attack 6 (unproved equivalence).} None introduced;
block-additivity is unconditional. \textbf{Pass.}

\paragraph{Attack 7 (compute-engine sync).} No engine.
\textbf{Pass.}

Cycle~1 closes with block-additivity as unconditional content.

\subsection{Cycle 2 (Per-Block Residual Isolation)}

\paragraph{Attack 1.} Add sources: Paskunas 2013 arXiv:1311.3048
(block-by-block projective-cover analysis); Colmez 2010
Ast\'erisque 330 (the Montreal functor bridging automorphic and
Galois sides). \textbf{Pass, extended disjoint set.}

\paragraph{Attack 2.} Paskunas's projective cover $P_{\rho}$ uses the
Montreal functor normalisation; this matches
\Cref{v4-arith:f1a3-BH:conj:BH-packaging}'s $P_{\tau}$ under the
BH-to-residual $\tau\mapsto\rho_{\tau}$ correspondence of
\Cref{v4-arith:f1a3-BH:prop:type-to-component-p2}. \textbf{Pass.}

\paragraph{Attack 3.} Block-local statements inside the ambient
abelian block category. \textbf{Pass.}

\paragraph{Attack 4.} The per-block hypothesis is named:
\Cref{v4-arith:w5-f:F1a3-AB:prop:A-residual-taxonomy}(2). No
further hypothesis needed for the structural factorisation.
\textbf{Pass.}

\paragraph{Attack 5.} Per-block identification is functorial in the
block; the product compatibility is handled by block-additivity.
\textbf{Pass.}

\paragraph{Attack 6.} (R.A) is the named per-block residual.
\textbf{Named.}

\paragraph{Attack 7.} No engine. \textbf{Pass.}

Cycle~2 closes: F1.a.3.A factors as unconditional block-additivity
plus per-block residual (R.A).

\subsection{Cycle 3 (Generic-Block $\mathrm{Ext}^{1}$-Vanishing)}

\paragraph{Attack 1.} Paskunas 2013 \S 6 (block centraliser regular
local of dim~$2$ or $3$); Auslander--Buchsbaum 1957 (Koszul
resolution on regular local); Paskunas 2015 arXiv:1308.3370
(Breuil--Mezard multiplicities). All disjoint from Bernstein--Deligne
and from Ben-Zvi--Nadler. \textbf{Pass.}

\paragraph{Attack 2.} Regular-local-ring convention uses the
reduced Koszul complex (Matsumura 1989 Ch.~6); matches
Gaitsgory--Rozenblyum Vol.~II for IndCoh on a regular local scheme.
\textbf{Pass.}

\paragraph{Attack 3.} Localisation at the generic point is
functorial on $A_{\rho}$-Mod; Morita equivalence is a morphism of
abelian categories, preserves Ext. \textbf{Pass.}

\paragraph{Attack 4.} The hypothesis ``$p\geq 5$, generic $\rho$,
fixed $\chi$'' is explicit in
\Cref{v4-arith:w5-f:F1a3-AB:prop:Ext1-generic-p5}. No implicit
hypothesis on the Langlands side (all inputs are from the block-algebra
structure). \textbf{Pass.}

\paragraph{Attack 5.} Generic-fibre Ext vanishing is a pointwise
statement on the generic fibre; no functoriality claim beyond
Morita. \textbf{Pass.}

\paragraph{Attack 6.} No new unproved equivalence.
\textbf{Pass.}

\paragraph{Attack 7.} No engine. \textbf{Pass.}

Cycle~3 closes: $\mathrm{Ext}^{1}$-vanishing at $p\geq 5$ on
generic locus is unconditional.

\subsection{Cycle 4 (Ordinary-Block Controlled Deformation)}

\paragraph{Attack 1.} Paskunas 2013 \S 10 (ordinary block);
Paskunas 2015 arXiv:1308.3370 (Breuil--Mezard multiplicity update);
Emerton--Paskunas arXiv:1509.04411 (reconciliation of earlier
versions). Disjoint from the generic-block source set.
\textbf{Pass.}

\paragraph{Attack 2.} Very-special sub-stratum convention follows
Paskunas 2013 Remark~10.5; matches Emerton--Paskunas's explicit
list. \textbf{Pass.}

\paragraph{Attack 3.} Ordinary block is a sub-block of the Bernstein
product with specific Frobenius-action signature on the
semisimplification of $\rho$. Same ambient as Cycle~1.
\textbf{Pass.}

\paragraph{Attack 4.} The ``very special'' sub-stratum is an explicit
enumerable case; its Breuil--Mezard multiplicity is either known
from Paskunas 2015 (non-very-special) or named as (R.B.1)
(very-special). Hypotheses named. \textbf{Pass.}

\paragraph{Attack 5.} Controlled-deformation statement is functorial
via the Paskunas framework. \textbf{Pass.}

\paragraph{Attack 6.} (R.B.1) is the named residual.
\textbf{Named.}

\paragraph{Attack 7.} No engine. \textbf{Pass.}

Cycle~4 closes: $\mathrm{Ext}^{1}$-vanishing on the ordinary block
reduces to (R.B.1), a strictly narrower very-special sub-stratum.

\subsection{Cycle 5 (Wild-Prime Extension)}

\paragraph{Attack 1.} Bushnell--Henniart 2006 \S 12--15 (wild types
at $p=2,3$); Auslander--Reiten--Smal\o{} 1995 (linear-quiver block
algebra theory). Paskunas 2013 does not cover $p\in\{2,3\}$
directly. \textbf{Pass.}

\paragraph{Attack 2.} Ext convention: the $A_{2}$-linear quiver
$L_{0}\leftarrow L_{1}\leftarrow L_{2}$ with
radical-square-zero relations matches
\Cref{v4-arith:bzn-rcompletion:prop:R4-Ext1-p3}. \textbf{Pass.}

\paragraph{Attack 3.} Block centraliser algebra at wild primes is
still complete Noetherian local (by general representation theory
of reductive groups over local fields). \textbf{Pass.}

\paragraph{Attack 4.} The key hypothesis is regularity of the
block centraliser at wild primes, named
\Cref{v4-arith:w5-f:F1a3-AB:prop:centraliser-wild}. Expected
dimensions are $2$ at $p=2$, $\leq 3$ at $p=3$, matching the
\Cref{v4-arith:bzn-rcompletion:prop:R4-Ext1-p2} and $p=3$
dimensions. \textbf{Pass, hypothesis explicit.}

\paragraph{Attack 5.} Regularity-to-Ext functoriality via Koszul
resolution. \textbf{Pass.}

\paragraph{Attack 6.} (R.B.2) is the named residual. \textbf{Named.}

\paragraph{Attack 7.} No engine. \textbf{Pass.}

Cycle~5 closes: $\mathrm{Ext}^{1}$-vanishing at $p\in\{2,3\}$
conditional on (R.B.2).

\section{Summary}
\label{v4-arith:w5-f:F1a3-AB:summary}

\begin{enumerate}
\item F1.a.3.A (BH packaging) factors into two parts: block-additivity
of Hochschild homology
(\Cref{v4-arith:w5-f:F1a3-AB:lem:HH-block-additive}, unconditional)
plus per-block identification
(\Cref{v4-arith:w5-f:F1a3-AB:prop:A-residual-taxonomy}(2), narrower
residual (R.A)).
\item F1.a.3.B (Paskunas wild-block Ext${}^{1}$) factors into three
regimes: generic blocks at $p\geq 5$ where
$\mathrm{Ext}^{1}=0$ on the generic locus unconditionally
(\Cref{v4-arith:w5-f:F1a3-AB:prop:Ext1-generic-p5}); the ordinary
block at $p\geq 5$ where $\mathrm{Ext}^{1}$ is controlled by
Paskunas 2015 and reduces to the very-special sub-stratum
residual (R.B.1); and the wild-prime extension at
$p\in\{2,3\}$ conditional on block centraliser regularity
(R.B.2).
\item The combined F1.a.3.A+B reduced ledger has four items:
(R.A), (R.B.1), (R.B.2), and the unchanged (R.C) of
Chapter~\ref{v4-arith:f1a3-BH}. Total paper-length
$13$--$26$ pp
(\Cref{v4-arith:w5-f:F1a3-AB:prop:F1a3-ledger-updated}), within the
$15$--$25$ pp estimate of Chapter~\ref{v4-arith:f1a3-BH}. The
partition is strictly finer: two original conjectures A and B become
three named sub-items, each strictly narrower.
\item The unconditional content extracted is: (i) block-additivity
of Hochschild homology on any orthogonal-idempotent block
decomposition
(\Cref{v4-arith:w5-f:F1a3-AB:lem:HH-block-additive}); (ii) residual
factorisation for F1.a.3.A
(\Cref{v4-arith:w5-f:F1a3-AB:prop:A-residual-taxonomy}); (iii)
generic-locus $\mathrm{Ext}^{1}$-vanishing at $p\geq 5$
(\Cref{v4-arith:w5-f:F1a3-AB:prop:Ext1-generic-p5} and
\Cref{v4-arith:w5-f:F1a3-AB:cor:Ext1-zero-generic-locus});
(iv) updated ledger
(\Cref{v4-arith:w5-f:F1a3-AB:prop:F1a3-ledger-updated}).
\item No $\mathrm{ProvedHere}$ inscription is attempted for the
full F1.a.3.A or F1.a.3.B statements; closure is reductive, not
unconditional. Attack--heal cycles pass all seven attacks at each
cycle; named residuals are three (R.A, R.B.1, R.B.2).
\end{enumerate}

The chapter converts two original conjectures into three strictly
narrower named conjectures plus four unconditional reductions.
Honesty over closure.

\section{Bibliography Stanza}
\label{v4-arith:w5-f:F1a3-AB:bibliography}

Primary sources used in this chapter (all external to the
arithmetic branch's inherited catalog):

\begin{itemize}
\item I.~Bernstein, ``Le centre de Bernstein'', in Deligne's notes
(1984).
\item J.~Bernstein and P.~Deligne, ``Le centre de Bernstein'', in
\emph{Repr\'esentations des groupes r\'eductifs sur un corps local}
(1984).
\item C.~J.~Bushnell and G.~Henniart, \emph{The Local Langlands
Conjecture for $GL(2)$}, Grundlehren der mathematischen Wissenschaften
335, Springer (2006).
\item P.~Colmez, ``Repr\'esentations de $GL_{2}(\mathbb{Q}_{p})$ et
$(\varphi,\Gamma)$-modules'', Ast\'erisque 330 (2010), 281--509.
\item V.~Paskunas, ``The image of Colmez's Montreal functor'',
Publ.~IHES 118 (2013), 1--191, arXiv:1311.3048.
\item V.~Paskunas, ``On the Breuil--M\'ezard conjecture'',
arXiv:1308.3370 (2015).
\item M.~Emerton and V.~Paskunas, ``On effective local Langlands
for $GL(2)/\mathbb{Q}_{p}$'', arXiv:1509.04411 (2015).
\item M.~Auslander, I.~Reiten, and S.~Smal\o{},
\emph{Representation Theory of Artin Algebras}, Cambridge Studies
in Advanced Mathematics 36, CUP (1995).
\item H.~Matsumura, \emph{Commutative Ring Theory}, Cambridge
Studies in Advanced Mathematics 8, CUP (1989).
\item D.~Ben-Zvi, J.~Francis, and D.~Nadler, ``Integral transforms
and Drinfeld centers in derived algebraic geometry'',
arXiv:0805.0157 (2008).
\item D.~Ben-Zvi and D.~Nadler, ``The character theory of a complex
group'', arXiv:0904.1247 (2009).
\item J.~Lurie, \emph{Higher Algebra}, manuscript (2017 version).
\end{itemize}

All sources are primary. No citation chain passes through
Chapter~\ref{v4-arith:f1a3-BH}'s conjectural apparatus for the
derivation of unconditional results. Disjointness audit:
\begin{itemize}
\item Block-additivity (\Cref{v4-arith:w5-f:F1a3-AB:lem:HH-block-additive})
uses only Lurie HA and Ben-Zvi--Francis--Nadler; disjoint from
Paskunas and BH sources.
\item Generic-locus $\mathrm{Ext}^{1}$-vanishing
(\Cref{v4-arith:w5-f:F1a3-AB:prop:Ext1-generic-p5}) uses only
Paskunas 2013 and Matsumura (Koszul resolution on regular local);
disjoint from Ben-Zvi--Nadler and Bernstein--Deligne.
\item Ordinary-block controlled deformation
(\Cref{v4-arith:w5-f:F1a3-AB:prop:Ext1-ordinary-controlled}) uses
Paskunas 2015 on top of Paskunas 2013; one citation chain, two
disjoint inputs.
\end{itemize}
Disjointness holds throughout; no tautological registration.
